\documentclass[../../main]{subfiles}

\begin{document}

\section{Estruturas com uma Operação}
\label{sec:matematica:algebra:estruturas-algebricas:uma-operacao}

Nesta seção consideramos estruturas algébricas sobre a assinatura
\[
	\Sigma = \{\ast : 2\},
\]
isto é, uma única operação binária.

Uma estrutura desse tipo consiste em um conjunto não vazio $A$ munido de uma interpretação
\[
	\ast^A : A \times A \to A.
\]

A partir dessa mesma assinatura, diferentes estruturas são obtidas pela adição de axiomas.

\subsection{Magma}

Um magma é um modelo da assinatura $\Sigma = \{\ast : 2\}$ sem axiomas adicionais.

Ou seja, qualquer par $(A,\ast)$ onde
\[
	\ast : A \times A \to A
\]
é uma operação binária.

Não são exigidas propriedades como associatividade, identidade ou inversos.

\subsection{Semigrupo}

Um semigrupo é um magma $(S,\ast)$ que satisfaz o axioma de associatividade:
\[
	\forall a,b,c \in S,\quad a \ast (b \ast c) = (a \ast b) \ast c.
\]

Esse axioma garante que a ordem de associação de operações não altera o resultado de expressões finitas.

\subsection{Monóide}

Um monóide é um semigrupo $(M,\ast)$ tal que existe um elemento $e \in M$ satisfazendo:
\[
	\forall a \in M,\quad a \ast e = e \ast a = a.
\]

O elemento $e$ é chamado elemento neutro.

\subsection{Grupo}

Um grupo é um monóide $(G,\ast)$ tal que:
\[
	\forall a \in G,\ \exists a^{-1} \in G \text{ tal que } a \ast a^{-1} = a^{-1} \ast a = e.
\]

O elemento $a^{-1}$ é chamado inverso de $a$.

\subsection{Grupo abeliano}

Um grupo $(G,\ast)$ é dito abeliano se adicionalmente satisfaz:
\[
	\forall a,b \in G,\quad a \ast b = b \ast a.
\]

\subsection{Hierarquia estrutural}

As estruturas definidas nesta seção formam uma cadeia de especialização:
\[
	\text{grupo} \subset \text{monóide} \subset \text{semigrupo} \subset \text{magma}.
\]

Cada nível é obtido pela adição de axiomas à mesma assinatura $\Sigma = \{\ast : 2\}$.

\subsection{Observação conceitual}

Essa abordagem evidencia que estruturas algébricas não diferem na natureza das operações, mas sim nos axiomas que essas operações satisfazem.

Isso será fundamental para a definição de morfismos e para a construção de estruturas mais complexas como anéis e corpos.

\end{document}
